\section{Related Work}
\label{sec:related}

\looseness -1 \textbf{Forecasting World Events.} Much prior work in Machine Learning and Statistics has focused on forecasting numeric or time-series data~\citep{BoxJenkins1976} in diverse domains like weather~\citep{Richardson1922}, econometrics~\citep{Tinbergen1939} or finance~\citep{Cowles1933}. Instead, our work focuses on the prediction of discrete world events, with both questions and answers described in natural language, also called \textit{judgemental forecasting}~\citep{tetlock2016superforecasting}. In the rest of our paper, we refer to this as \textit{forecasting} for brevity. In prior work on evaluating language models for forecasting~\citep{zou2022forecasting, karger2024forecastbench}, questions are primarily sourced from prediction markets like Metaculus, Manifold, and Polymarket. Prediction markets provide a platform for online participants to register predictions on questions like ``Will Donald Trump win the US Presidential Election in 2024?", which mostly have binary, yes or no, outcomes and have rapidly grown in popularity over the last few years.  \vspace{0.15cm}

\looseness -1 \textbf{Evaluating LLMs for Forecasting.} New information (before the event resolves) benefits forecasting. Thus, LLM forecasting work~\citep{zou2022forecasting, halawi2024approaching} provides relevant retrieved articles to models~\citep{lewis2020retrieval} often obtained via web-search APIs. \citet{paleka2025pitfalls} discuss pitfalls of LLM forecasting evaluations, including leakage of outcomes from online search in backtests, and distributional biases of prediction market questions. To avoid these issues, we use static, monthly snapshots of global news for retrieval and creating questions. \citet{jin-etal-2021-forecastqa} ask humans to create forecasting questions, while \citet{dai2024llms} try to automate this process with LLMs. However, their questions pre-define a few outcomes to choose from. ~\citet{guan2024openep, wang2025openforecast} evaluate models on open-ended forecasts, but we go a step further by showing how to train models for this task. \vspace{0.15cm} 

\textbf{Reinforcement Learning for LLMs.} \citet{shao2024deepseekmath} proposed \textit{Group Relative Policy Optimization} (GRPO), an RL algorithm that only uses outcome rewards. This approach has been highly successful in training LLMs to \textit{reason} about well-specified coding~\citep{jain2024livecodebench} and exam-style questions across domains~\citep{phan2025humanity}. Instead, forecasting requires LLMs to reason about uncertainty. \citet{halawi2024approaching} proposed training language models for forecasting by Supervised Finetuning (SFT) on chain of thought traces that lead to brier scores better than the prediction market aggregate. In the same setting of binary forecasting questions, \citet{turtel2025llms} optimize brier scores using GRPO, while ~\citet{damani2025beyond} extend it to short answer questions in other domains. We depart from these works in showing how to synthesise large-scale forecasting training data from daily news, to train models at open-ended reasoning about the future. 


\section{Open-Ended Forecasting}
\label{sec:open-ended}

\textbf{Motivation.} The forecasting task we study is \textit{open-ended} in two key ways: 
\begin{enumerate}
    \item It allows expressing arbitrary natural language forecasting questions
    \item It may not have a structured outcome set, unlike numeric or categorical predictions. This differentiates it from both time-series forecasting, and prediction markets.
\end{enumerate}
For example, prediction markets are dominated by binary (yes/no) or multiple choice questions. While this design is easy to score, the most foresight often lies in predicting the unexpected, or when a large number of possibilities could occur. The most important questions to forecast---such as scientific breakthroughs, geopolitical shocks, or technological disruptions---often emerge as \emph{unknown unknowns}: possibilities not anticipated, and hard to enumerate. Thus, in this work, we focus on training models to make open-ended predictions like "Which company will the US Government buy a >5\% stake in by September 2025?". Such questions require exploration and imagination, rewarding novel hypotheses that turn out to be correct, rather than just distributing probabilities over a known set of outcomes.\vspace{0.15cm} 

\textbf{Background.} LLM weights are frozen after training, especially for open-weight models. Any event that happens after the last date in the training corpus is in the future for the LLM. This provides a time window to collect questions for training models to reason about future events. Similarly, their evaluation involves testing on questions resolving after the cutoff date of the training data, called \textit{backtesting}~\citep{Tashman2000}. While prior work has relied on prediction market questions as training data, this has three key problems:
\begin{enumerate}
    \item The questions are created by humans, which makes them low in number~\citep{paleka2025pitfalls}. This becomes a bottleneck for scaling training data, which has been an essential component in the success of LLMs~\citep{kaplan2020scalinglawsneurallanguage, lu2025internet}.
    \item Most questions have binary outcomes, which creates a 50\% baseline success rate. This leads to noisy rewards in outcome-based RL, which means even incorrect reasoning has a high chance of being reinforced.
    \item Each platform has a skewed distribution of events. All overrepresent US political news, along with their specific focus such as crypto-currency price movements in Polymarket, technology in Metaculus, personal life of users in Manifold, and sports events on Kalshi~\citep{paleka2025pitfalls}. 
\end{enumerate}   
These limitations motivate us to explore alternate ways to create forecasting questions.\vspace{0.15cm} 

\looseness -1 \textbf{Setup.} Let $\mathcal{X}$ be the set of open-ended forecasting questions; and $\mathcal{Y}$ the set of short textual answers. We provide a language model $\pi_\theta$ a question $x \in \mathcal{X}$, for which we already know the ground-truth outcome $y^\star$ as it has resolved in the real-world. We ask the model to report its best guess prediction $y$, and the probability $q$ of it being the true outcome. \vspace{0.15cm} 

\textbf{Measuring Accuracy.} We measure accuracy by checking if the model's attempted answer $y$ matches with the ground truth outcome $y^\star$, using another language model to test for semantic equivalence (for example ``Geoffrey Hinton'' = ``Geoffrey Everest Hinton'') consistent with recent frontier benchmarks~\citep{wei2024measuringshortformfactualitylarge, phan2025humanity}. For evaluations, we use Llama-4-Scout~\citep{meta_llama4_2025}, as in a recent study~\citep{chandak2025answer}, it at matching answers, it has inter-human levels of alignment with human judgments. During training, we use Qwen3-4B in non-thinking mode, as it achieves high alignment levels for its size in the same study. We find the two models agree on $\sim 97\%$ grading responses, and manual validation ensures they are accurate in $\ge 95\%$ cases.\vspace{0.15cm} 

\textbf{Measuring Calibration.} We adapt the multi-class Brier scoring rule \citep{mucsanyi2023trustworthy} for free-form responses as follows (details in \Cref{sec:freeform_brier}):
$$
S^\prime(q, y, y^*) = \begin{cases}
    1 -(q - 1)^2, & \text{if } y \equiv y^*\\
     - q^2, & \text{if } y \neq y^*
    \end{cases} \\
    \label{eq:rlcr}
$$
\looseness -1 \textbf{Interpretation.} Predicting an event with a probability $q=0$ returns a baseline score of $0$ regardless of the prediction $y$. Correct predictions receive positive scores while incorrect predictions negative. For brevity, we call $S^\prime(q, y, y^*)$ \emph{Brier score} throughout this paper. Our Brier score is equivalent to the reward metric used by~\citet{damani2025beyond}. They show this is a proper scoring rule, incentivizing both high accuracy and truthful reporting of probability on the answer that seems most likely. For completeness, we discuss this further in \Cref{sec:freeform_brier}.\vspace{0.15cm} 

\textbf{Training Algorithm: GRPO~\citep{shao2024deepseekmath}.} 
We train LLMs using outcome-based reinforcement learning on our dataset. For each prompt $x$, we draw $K$ completions $\{(y_i,p_i)\}_{i=1}^K \sim \pi_\theta(\cdot\mid x)$ and compute rewards $r_i=R(y_i,p_i;y^\star)$. However, following prior work \citep{damani2025beyond, turtel2025outcome}, we do not divide by \emph{remove} the group standard deviation when computing advantages, as this stabilizes updates in settings like ours where reward variance can be small. \vspace{0.15cm} 

\textbf{Initial Policy: Qwen3 Thinking~\citep{yang2025qwen3technicalreport}.} We start with the 8B thinking model. For Qwen3 models, no official knowledge–cutoff date is reported. When queried directly, the models return inconsistent cutoff dates (most often \emph{October 2023} or \emph{June 2024}). Usually, they treat questions about 2024 as being in the future. Since the model weights were released and frozen in April 2025, we train up to this date, and use the period between May to August 2025 for testing. In the Appendix, we show large improvements from our training on even the Llama and Gemma models in~\Cref{sec:extra_forecasting_results}.
